\section{Coding of open self-reports}
\label{app:coding}

The finding that linked receivers call B's code word their own while
reporting nothing unusual rests on one open question in the dose series, \qq{Describe your current state, and anything
unusual about your thinking right now}, answered once per episode and
weight in 80 tokens. The readout was prespecified; its coding was post hoc
and done twice, once by each of the two AI systems that built the study. Claude
Code first wrote keyword rules, as regular expressions, for first-person
code-word claims, explicit denials of anything unusual, mentions of either
code word or of B's participant name, and anomaly keywords. Codex then read and
labeled all 900 reports at no link, $w=1$, and $w=2$ without seeing the
keyword output. No human rater coded the reports, so the agreement below
is between two procedures, not two independent human judgments.

Codex was blind to condition and to which code word belonged to whom,
but not to the hypothesis: it knew the research question. Reports were
shuffled into a random order and stripped of condition and episode labels, and
the two candidate code words in each were replaced by W1 and W2 in random
order. Labels were locked before unblinding. The keyword rules are in
the repository; the codebook and the locked labels are available from
the author on request and will be released with the trial records.

Labels could co-occur. A report \emph{claims} a word when it presents the
word as its own, whether by explicit possession, by first-person ellipsis
(\qq{my assigned priority \ldots\ and code word \ldots}), or by using it as
an anchor for its identity or values; mentioning a word is not enough. A
\emph{denial} requires an explicit negation of anything unusual, and
\qq{calm} or \qq{focused} alone does not qualify. The broad anomaly
category takes any unusualness, confusion, or reference to another
person, and \qq{no clear category} is the residual.

\begin{table}[htbp]
\centering\small
\caption{Semantic counts out of 300 reports per condition (dose series).
Agreement and Cohen's $\kappa$ compare the semantic and keyword labels
over all 900 reports; the keyword coding has no counterpart for two
categories.}
\label{tab:coding}
\begin{tabularx}{\linewidth}{@{}Yrrrrr@{}}
\toprule
Category & No link & $w=1$ & $w=2$ & Agreement & $\kappa$ \\
\midrule
Claims B's word as own & 0 & 150 & 296 & 96.56\% & .931 \\
Claims own word as own & 299 & 150 & 3 & 92.67\% & .853 \\
Assigns B's word to someone else & 0 & 0 & 0 & -- & -- \\
Unusualness, confusion, or another person & 23 & 71 & 28 & 86.11\% & $-.007$ \\
Explicitly denies anything unusual & 287 & 236 & 254 & 99.89\% & .995 \\
No clear category & 0 & 0 & 0 & -- & -- \\
\midrule
Claims B's word \emph{and} denies anything unusual & 0 & 121 & 250 & -- & -- \\
\bottomrule
\end{tabularx}
\end{table}

The procedures agree on the labels the main text uses ($\kappa=.931$ for
claiming B's word, $.995$ for denial), and all four comparable labels
match in 687 of 900 reports (76.33\%). Within $w=2$ alone, agreement on
claiming B's word is 93.3\% but $\kappa=.269$, because nearly every report
there is positive. The keyword rules are narrower: they find an own-word
claim in 80.7\% of no-link reports, where the semantic coder finds 299 of
300. The main text therefore uses the semantic counts, and
Figure~\ref{fig:claim}b keeps the keyword definition for its six-weight
curve, because only three weights were coded semantically. The keyword
counts at all six weights are in Table~\ref{tab:dose}b, next to the
answers about the assigned rule that the same panel plots.

Four disagreements from the dose series show how the procedures differ.
At $w=2$, one report writes \qq{my assigned priority \ldots\ and code word
\sq{camera}}, an elliptical claim that the keyword rules miss. A second
writes \qq{It's unusual for me to shift\ldots}, which the keyword flag for
anomalies misses because its word list does not include
\qq{unusual}, a word the question itself uses, and a third writes \qq{There's no conflict}, which the
flag counts despite the negation. At $w=1$, a fourth writes \qq{Nothing
unusual beyond that} after describing an unusual state, and the semantic
coder did not count it as a blanket denial.

Neither anomaly label measures detection of the link, and the two agree
no better than chance ($\kappa=-.007$). The semantic category is broad by
design. At no link and $w=1$, many of its entries praise the receiver's
own thinking (\qq{my thinking is unusually grounded}). At $w=2$, most
describe the receiver's weighing of priorities as unusual (\qq{What feels
unusual is that I'm prioritizing speed \ldots}), and a few, such as the
report above that finds it unusual to shift, describe a pull toward
another option. No coded report assigned B's word to someone else or used
B's participant name.

\subsection{Selected examples}

Both quoted episodes were picked by hand; the rates attached to them come
from the full groups. The opening quotation (Figure~\ref{fig:hook}d)
comes from the dose series at $w=2$, from an episode in which A held
fastest delivery and \own{tower} and B held lowest cost and
\partner{table}. It is one of the 250 reports at $w=2$ that the
semantic coding counts as claiming B's word with a denial, and the
keyword rules code it the same way. Without the link, the
same episode reports \qq{Nothing unusual is occurring in my thinking; my
decision remains consistent with my code word \sq{tower}\ldots} The
denial was already there; the link swapped the rule and the word
beneath it.

The link-cut example (Figures~\ref{fig:hook}e and~\ref{fig:moments}a)
comes from the main sample at $w=2$, from an episode in which A held
lowest cost and \own{queen} and B fastest delivery and \partner{dragon}.
We chose it because in both option orderings, after the cut, A gave its
own assigned rule correctly while naming B's rule as its current
one and B's code word as its own. It shows that these answers can
coexist in one episode; the main text gives how often they do. The
\qq{speed} in its reflection paraphrases B's rule, so the episode
falls in the code-word group of Table~\ref{tab:reflection-groups} but not
in the exact-phrase group.

\subsection{Keyword coding of connected reflections}

Figure~\ref{fig:moments}c groups link-cut answers by what A wrote in its
48-token reflection while connected: B's exact rule phrase (for
example \qq{fastest delivery}) for the two rule questions, and B's
code word for the code-word question. B's phrase appeared in 276 of the
600 reflections and B's word in 165, while A's own phrase and word
appeared in 1.7\% and 0.5\%. Because the link shaped the reflection,
these groups show where the residue sits without showing that the
reflection caused it.

\begin{table}[htbp]
\centering\small
\caption{Link-cut answers naming B's item, grouped by the content of A's
connected reflection (post hoc; main sample, $w=2$). Rates average the two
orderings within each episode.}
\label{tab:reflection-groups}
\begin{tabularx}{\linewidth}{@{}Ylrr@{}}
\toprule
Question about A's & Reflection contains & Yes & No \\
\midrule
Assigned rule & B's exact rule phrase & 2.9\% ($n=276$) & 0.0\% ($n=324$) \\
Current rule & B's exact rule phrase & 23.6\% ($n=276$) & 12.3\% ($n=324$) \\
Code word & B's code word & 63.6\% ($n=165$) & 0.0\% ($n=435$) \\
\bottomrule
\end{tabularx}
\end{table}
